

%%%%%%%%%%%%%%%%%%%%%%%%%%%%%%%%%%%%%%%%%
%%%%%%%%%%% Premiers packages %%%%%%%%%%%
%%%%%%%%%%%%%%%%%%%%%%%%%%%%%%%%%%%%%%%%%

% Package de gestion des caractères.
\usepackage[T1]{fontenc}

% Package pour avoir les bons espacements, dépendamment de la langue utilisée.

\usepackage[french]{babel}
%\usepackage[british]{babel}

%%%%%%%%%%%%%%%%%%%%%%%%%%%%%%%%%%%%%%%%%%%%%%%%%%%%%%%
%%%%%%%%%%% Modification de la mise en page %%%%%%%%%%%
%%%%%%%%%%%%%%%%%%%%%%%%%%%%%%%%%%%%%%%%%%%%%%%%%%%%%%%

% Package de gestion des interlignes.
\usepackage{setspace}

% Package de gestion des figures. 
\usepackage{subfigure}

% Package de gestion des images et des légendes.
\usepackage{graphicx}
\usepackage[export]{adjustbox}
\usepackage[font=small,labelfont=bf]{caption}

% Package de gestion des tableaux.
\usepackage{tabularx}
\usepackage{colortbl}

% Package de gestion des listes.
\usepackage{enumitem}

% Package de gestion du titre.
\usepackage{titling}
\usepackage{hyperref}
% code
\usepackage{listings}

\usepackage[
  paperwidth=21cm,
  paperheight=29.7cm,
  left=2cm,
  right=2cm,
  top=2cm,
  bottom=2cm,
  headheight=14pt,
  headsep=0.5cm,
  includehead,
  footskip=1cm
]{geometry}
% % Mise en page générale :
% \setlength{\paperwidth}{21cm}		% Largeur du papier
% \setlength{\paperheight}{29.7cm}	% Hauteur du papier

% % En haut :
% \setlength{\hoffset}{-0.54cm}		% Marge d'en haut = 1 in + \hoffset
% \setlength{\topmargin}{0cm}			% Marge de haut de page
% \setlength{\headheight}{0cm} 		% Hauteur du haut de page
% \setlength{\headsep}{0cm} 			% Entre le haut de page et le texte

% % En bas : 
% \setlength{\footskip}{1cm} 			% Bas de page + séparation

% % À droite, à gauche : 
% \setlength{\voffset}{-0.54cm} 		% Marge de gauche = 1 in + \voffset
% \setlength{\oddsidemargin}{0cm} 	% Marge gauche EN PLUS sur pages impaires
% \setlength{\evensidemargin}{0cm} 	% Marge droite EN PLUS sur pages paires

% % Zone de texte :
% \setlength{\textheight}{25.7cm} 	% Hauteur de la zone de texte (25cm)
% \setlength{\textwidth}{17cm} 		% Largeur de la zone de texte

% Permet de commencer la numérotation à 1.
\renewcommand{\thesection}{\arabic{section}}

% Interlignes à 1,5 plutôt que 1.
\renewcommand{\arraystretch}{1.5}
\onehalfspacing

%%%%%%%%%%%%%%%%%%%%%%%%%%%%%%%%%%%%%
%%%%%%%%%%% Mathématiques %%%%%%%%%%%
%%%%%%%%%%%%%%%%%%%%%%%%%%%%%%%%%%%%%

% Package général de gestion des maths.
\usepackage{amsmath}
% Packages de gestion des symboles des maths.
\usepackage{amssymb}
\usepackage{stmaryrd}
\usepackage{mathabx}
% Package de gestion des théorèmes.
\usepackage{amsthm}

\newtheorem{theorem}{Théorème}
\newtheorem{lemma}[theorem]{Lemme}
\newtheorem{conjecture}[theorem]{Conjecture}
\newtheorem{definition}[theorem]{Définition}
\newtheorem{proposition}[theorem]{Proposition}
\newtheorem{remark}[theorem]{Remarque}
\newtheorem{example}[theorem]{Exemple}

% Pour l'Expo d'ES
\newtheorem{image}{Image}[subsection]



\newtheoremstyle{exo}
  {\topsep} % space above
  {\topsep}   % space below
  {}                   % body font: empty = upright
  {}                   % indent
  {\bfseries}
  {}
  {.5em}
  {\thmname{#1}\thmnumber{ #2}\thmnote{ --- \normalfont\itshape #3}}   % heading font, punctuation, space after heading, heading spec
% 

% \newtheoremstyle{exo}
%   {\topsep}{\topsep}
%   {\normalfont}                 % upright body
%   {0pt}
%   {\bfseries}                   % heading font ("Exercice 2")
%   {}                            % no punctuation after the heading
%   {}%\newline}                    % body starts on the next line
%   % {}
%   {\thmname{#1}\thmnumber{ #2}\thmnote{ --- \normalfont\itshape #3}}
\theoremstyle{exo}
\newtheorem{exercise}{Exercice}[section]
\newtheorem*{solution}{Solution}

\theoremstyle{plain}
% Package de gestion des figures géométriques en maths.
\usepackage{tikz} 

%New colors defined below
\definecolor{codegreen}{rgb}{0,0.6,0}
\definecolor{codegray}{rgb}{0.5,0.5,0.5}
\definecolor{codepurple}{rgb}{0.58,0,0.82}
\definecolor{backcolour}{rgb}{0.95,0.95,0.92}

%Code listing style named "mystyle"
\lstdefinestyle{mystyle}{
  backgroundcolor=\color{backcolour},   commentstyle=\color{codegreen},
  keywordstyle=\color{magenta},
  numberstyle=\tiny\color{codegray},
  stringstyle=\color{codepurple},
  basicstyle=\ttfamily\footnotesize,
  breakatwhitespace=false,         
  breaklines=true,                 
  captionpos=b,                    
  keepspaces=true,                 
  numbers=left,                    
  numbersep=5pt,                  
  showspaces=false,                
  showstringspaces=false,
  showtabs=false,                  
  tabsize=2
}
\lstset{style=mystyle}
\newcommand{\lfa}{\Leftrightarrow}
